(1) 以桌面高度为零势面，悬挂段质量 \(\lambda y\)，其质心在 \(y/2\) 的下方，故全链势能 \(U(y)=-\lambda gy^2/2\)。整链同速，动能 \(K=\lambda Lv^2/2\)。理想无损条件下 \(K+U=-\lambda gy_0^2/2\)，即 \(Lv^2=g(y^2-y_0^2)\)。该结论仅用于经过圆角的连续可弯链，不能直接用于存在非弹性折角撞击的真实链。

(2) 对能量式求导，在 \(v>0\) 时有 \(2Lv\dot v=2gyv\)，延续到初时给 \(\ddot y=(g/L)y\)。记 \(\omega=\sqrt{g/L}\)，满足 \(y(0)=y_0,\dot y(0)=0\) 的解为 \(y(t)=y_0\cosh(\omega t)\)，\(v(t)=y_0\omega\sinh(\omega t)\)。直接代回能量恒等式，利用 \(\cosh^2-\sinh^2=1\)，可核验解的正确性。

(3) 完整离桌条件 \(y(t_*)=L\)，因此 \(t_*=\sqrt{L/g}\operatorname{arcosh}(L/y_0)\)，只在链尾尚未离桌之前使用上述解。相应速率为 \(v_*=\sqrt{g(L^2-y_0^2)/L}\)，与能量式直接代入 \(y=L\) 相符。\(y_0\to L^-\) 则 \(t_*\to0\) 且 \(v_*\to0\)；\(y_0\to0^+\) 的启动时间趋无穷，因为模型在严格零悬垂长度无初始驱动力。

(4) 若仅考察悬挂质量 \(\lambda y\)，其动量变化为 \(d(\lambda yv)/dt=\lambda y\dot v+\lambda v^2\)，其中额外项来自新增链节携带的向下动量。开口控制体动量方程必须在边界计入 \(\lambda v^2\) 的通量及圆角区张力贡献；仅写 \(F=(\lambda y)\dot v\) 并把它当作全系统力平衡会忽略质量交换效应。对整链使用能量法绕开了对局部几何的虚假简化。
